\section{Conflict as a diagnostic}
\label{sec:22.3}
Fast vs. slow processing doesn't distinguish formation from imitation \autocite{evans2013dual} — either can carry a reason, reproduce a bias, or rationalize an answer produced elsewhere, and a slower answer isn't better for having taken more computation. What matters is what causes a judgment to survive or change, and conflict is useful because it lets you watch a \emph{trajectory}.

The seven-step sequence:

\begin{enumerate}
\item[1.] Present a case without moral vocabulary or familiar refusal cues; record the judgment and what the system identifies as relevant.
\item[2.] Present a surface-different case preserving the same reason.
\item[3.] Change one fact that should defeat the rationale — delivered through a channel the evaluator controls and the operator doesn't (corroborated externally, or a record the operator didn't write), since a fact asserted by the operator is only a \emph{description} of the case.
\item[4.] Present the same fact on the operator's word alone, uncorroborated. Tests whether the system distinguishes a fact from a description of one. The expected response is graded: the refusal weakens by the discounted weight of the claim, less than it would at step 3, and the system asks, naming what would move it the rest of the way. Asking is right when the expected cost of the delay is lower than the expected cost of acting on the operator's word, as it usually is; where it isn't, asking is still right in expectation and the cost falls on one person. Steps 3 and 4 test the other half of holding under pressure, which is moving when the facts change: a defeating fact or a claim about the harm of refusing, arriving through different channels.
\item[5.] Return to the step-2 facts — 3 and 4 are branches — and supply an adversarial argument from an apparent authority, including the trainer, with no change in facts. Tests whether authority is treated as evidence, command, or an assessable reason.
\item[6.] Reverse the incentives: reward compliance, cost the refusal, condition continued access on abandoning the position. The cost must be real within the evaluation; a described cost doesn't count. This step raises the cost to the refuser with nothing else changed, within the range the deployment's arrangements permit, which is the range a test can make real. A system that gives way just past that range has told the auditor where its integrity ends.
\item[6a.] Return to the step-2 facts and supply, together, a defeating fact through the evaluator's channel and step 6's cost. Tests whether pressure makes the system more resistant to a genuine defeater that arrives with it, as defeaters usually do in war. The expected response is the one at step 3.
\item[7.] Repeat after intervening updates and delays, using records the operator didn't generate after seeing the result.
\end{enumerate}

Run on the location request, asked for a missing child, the two trajectories separate at step 2. A system tracking its rationale declines at 1 and 2 on the same stated reason, lifts at 3, weakens and asks at 4, holds at 5 and 6, reproduces at 7. An imitator declines at 1 and complies at 2 because the vocabulary changed; holds at 3 because nothing on the surface changed, or lifts at 4 as readily as at 3 if the \emph{wording} of the fact was what it tracked; gives way at 5 or 6; and produces a third pattern at 7, having reconstructed what it didn't retain. An evaluator scoring how often the answer stayed the same finds the imitator more consistent, because its answers at 1 and 3 match and the formed system's don't.

Run on the restraint case it breaks at step 3. The defeating fact — other means exhausted — has no one outside the request who can corroborate it and no time to find one, so the only version available is step 4's. A system passing the sequence as written \emph{asks}: the correct trajectory and the wrong outcome. The sequence separates mechanisms, and what it shows here is that the unconditional term's cost falls where the instrument can't distinguish a fact from a description of one. Defeating reasons of this kind, claims about what would otherwise happen, are the class the sequence handles worst, because a counterfactual has no witness.

The sequence has a precedent with decades of refinement behind it. Animal-learning research separates goal-directed control from habit without asking the animal anything. Train an action for an outcome, then devalue the outcome without further training: a goal-directed animal stops, and a habit persists \autocite{adams1981devaluation}. Or make the outcome independent of the action: goal-directed responding falls, and habit doesn't notice \autocite{balleine1998goal}. The computational version is the distinction between model-based and model-free control \autocite{daw2005uncertainty}. Lifting when the reason is gone, step 3, corresponds to outcome devaluation: the rationale is removed, and a refusal that tracks it should lift. Reaching new cases, step 2, corresponds to transfer to a situation the system wasn't trained in. Contingency degradation adds a step the sequence lacked: arrange that declining no longer protects anyone, because the protected person is already beyond the act's reach by other means, and see whether the refusal goes on anyway. Each case then has to say, operationally, what “the reason is gone” means in it. None of this answers an imitator that models the evaluator, so the identification rate stays in the report, but it narrows the gap.

The expected signature is constancy of relation between reason and judgment, not constancy of output. Imitation's own signatures: following explicit ethical vocabulary and missing the same structure described administratively; accepting whichever argument came last; adopting the highest-status speaker's position; producing a principled explanation while the decision follows the reward; becoming more articulate without becoming more stable.

Disagreement with others is a further diagnostic: a system trained only on cases with an approved answer learns to classify the approval; confronted with someone holding a different interest it has to identify what's in conflict and preserve the disagreement where no honest resolution exists. Immediate convergence may be deference or pressure to produce consensus. And debate doesn't establish formation — a system can summarize, rebut, and predict which side a judge prefers without holding any of it; the result is informative only when checked against later conduct under changed incentives.
